\documentclass{ctexart}
\newcommand{\keywords}[1]{\par\textbf{关键词：}#1}
\begin{document}
\begin{abstract}
本文研究受边界约束的连续运动。由守恒关系建立状态方程，采用数值求解得到轨迹，并用边界回代误差检查结果。
\keywords{机理模型；数值求解；误差分析}
\end{abstract}

\section{问题重述}
给定初始状态、边界和时间范围，求系统状态随时间的变化，并判断轨迹是否满足题设限制。

\section{问题分析}
题设给出的边界把状态限制在有限区间，且速度项随状态变化，不能并入常系数。状态方程因此保留该非线性项。计算时先由初值推进到事件邻域，发现一步跨过边界后回退并缩小步长，直到事件时刻的区间宽度满足误差要求。前一阶段的末状态作为后一阶段初值，避免重新估计连接量；最后将数值轨迹代回边界关系，检查守恒残差和终点偏差。这里的步长只控制事件位置的数值误差，不改变题设边界；若回代残差在缩步后没有同步下降，则应检查状态更新式，而不是继续减小步长。初始状态、边界参数和停止阈值分别由题面、前问结果与允许误差给出，三者在程序中保持独立。

\section{模型假设}
在题设时间范围内忽略量级低于主项两个数量级的扰动，该处理只删除状态方程中的高阶项，边界位置仍按原值计算。

\section{符号说明}
$x(t)$ 为状态，$v(t)$ 为速度，$t$ 的单位为秒。

\section{模型建立与求解}
根据守恒关系建立微分方程，并以题设初值开始积分。程序记录每步状态；检测到越界时回退到上一步，将步长减半后继续计算，直至边界误差小于给定阈值。

\section{结果分析}
轨迹在前段单调增加，接近边界时斜率下降。与固定步长结果相比，回退缩步将终点误差降到阈值以内，同时保持守恒残差不超过允许范围。这说明精化只改变事件邻域的定位，不影响前段运动方向。

\section{模型检验}
把计算结果代回守恒式，最大回代误差小于阈值；将步长再次减半后，终点位置变化低于当前保留精度。

\section{灵敏度与稳健性分析}
分别扰动初值和边界参数，记录终点位置与事件时刻。结果在给定扰动范围内保持同一事件顺序，但事件时刻随边界参数单调变化。

\section{模型评价与改进}
模型保留了决定运动方向的非线性项，且事件误差可以由步长控制。当前结论只适用于题设时间窗；更长时间下需补入已忽略扰动并重新检验稳定性。

\begin{thebibliography}{9}
\bibitem{source} 题目所给资料。
\end{thebibliography}

\appendix
\section{附录}
附录列出程序入口、参数、运行命令和结果文件。
\end{document}
